% TOP_PROBLEM: {"id": 10000039, "title": "Rough-isometry invariance of percolation nonuniqueness", "category_id": 19, "category": {"id": 19, "name": "probability", "display_name": "Probability", "description": "Problems involving probability theory, stochastic processes, and random structures.", "slug": "probability", "order_index": 19, "created_at": "2026-07-31T15:26:25.671Z"}, "status": "partially_solved", "problem_number": "AMR-099-0039", "set_id": 13}
% FIRST_POSTED: 2026-10-01
% ENTRY_KIND: statement_only
% UnsolvedMath ID 10000039; source catalogue code AMR-099-0039
% !TeX program = lualatex
% Definition and sources only; this file is not an LLM solution attempt.
\documentclass[11pt,a4paper]{article}
\usepackage[margin=25mm]{geometry}
\usepackage{fontspec}
\setmainfont{lmroman10-regular.otf}[BoldFont=lmroman10-bold.otf,
 ItalicFont=lmroman10-italic.otf,BoldItalicFont=lmroman10-bolditalic.otf]
\usepackage{amsmath,amssymb,mathtools}
\usepackage{unicode-math}
\setmathfont{latinmodern-math.otf}
\AtBeginDocument{\let\setminus\smallsetminus}
\usepackage{xurl,hyperref}
\hypersetup{colorlinks=true,urlcolor=blue}
\setcounter{secnumdepth}{0}
\setlength{\parindent}{0pt}
\setlength{\parskip}{5pt}
\setlength{\emergencystretch}{3em}
\begin{document}
\section*{Rough-isometry invariance of percolation nonuniqueness}\label{10000039}
{\small UnsolvedMath ID 10000039 · AMR-099-0039 · Probability · AMR Open Problem Lists}
\subsection{Definitions and mathematical statement}
For an infinite connected bounded-degree graph $G$, define Bernoulli bond percolation by independently retaining edges with probability $p$. Let $N_\infty$ count its infinite open components, and set
\[
p_c(G)=\inf\{p:\mathbb P_p(N_\infty\ge1)>0\},\qquad
p_u(G)=\inf\{p:\mathbb P_p(N_\infty=1)=1\}.
\]
Both infima are over $[0,1]$; $p=1$ supplies the unique infinite component.

Suppose $G$ and $H$ are infinite connected graphs of bounded degree, and there is a map $f:V(G)\to V(H)$ with constants $A\ge1$, $B\ge0$ satisfying
\[
A^{-1}d_G(x,y)-B\le d_H(f(x),f(y))\le A d_G(x,y)+B
\]
for all $x,y$, with every vertex of $H$ at distance at most $B$ from $f(V(G))$. Prove or disprove
\[
p_c(G)<p_u(G)\quad\Longleftrightarrow\quad p_c(H)<p_u(H).
\]
The rough-isometry terminology here includes multiplicative distortion; the map need not preserve edges, degrees, or a transitive group action. Neither graph is assumed transitive.

On a general graph the set of parameters with a unique infinite cluster need not itself be an interval. Consequently the proposed invariant is precisely the inequality between the two displayed infima, not an unstated monotonicity assertion or a characterization of every intermediate parameter. A counterexample must give a genuine rough-isometry pair and verify opposite truth values of this numerical inequality. The statement concerns invariance of the strict threshold separation, not equality of either individual threshold across the pair.
\subsection{Short English statement}
Is strict separation of percolation and uniqueness thresholds preserved by coarse geometric equivalence of bounded-degree graphs?
\subsection{Sources}
\begin{itemize}
\item[S1] ULAM.AI and the UnsolvedMath Contributors, supplied problems.json, record 10000039 (AMR-099-0039), AMR Open Problem Lists. Catalogue identity and source transcription; CC BY 4.0.
\url{https://huggingface.co/datasets/ulamai/UnsolvedMath}.
\item[S2] Benjamini - Coarse Geometry and Randomness (Saint-Flour notes), Conjecture 9.11, PDF page 65. Original source indexed by AMR-099-0039.
\url{https://www.wisdom.weizmann.ac.il/~itai/stflouraug24.pdf#page=65}.
\item[S3] I. Benjamini, Coarse Geometry and Randomness, primary Saint-Flour notes; the numbered question and its surrounding definitions.
\url{https://arquivo.pt/noFrame/replay/20201231041548id_/http://www.wisdom.weizmann.ac.il/~itai/stflouraug24.pdf}.
\end{itemize}
\end{document}
